\begin{pastexam}[label={px:nsysu-114-dmalg-04}]{中山資工所資安組 114 學年度 第 5 題}
（計算）

{[10\%]} Find the generating function for the number of integer solutions to the equation
$c_1 + c_2 + c_3 + c_4 = 23$, where $-2 \le c_1, -4 \le c_2, -3 \le c_3 \le 6, 0 \le c_4$.
\end{pastexam}

\begin{pastsoln}
答案：所求的個數是
\[
f(x) = \frac{x^{-9}\,(1 - x^{10})}{(1 - x)^4}
\]
裡 $x^{23}$ 的係數，等價地是 $g(x) = \dfrac{1 - x^{10}}{(1 - x)^4}$ 裡 $x^{32}$ 的係數。 這個係數是
$\binom{35}{3} - \binom{25}{3} = 4245$。

\textbf{每個變數一個因式}。 生成函數的做法（資結書第~2.6~節）：變數 $c_i$ 能取的每個值 $v$
對應一項 $x^v$，四個因式相乘後，$x^{23}$ 的係數就是「從每個因式各挑一項、指數和為 $23$」的
方法數，也就是解的個數。
\begin{align*}
c_1 \ge -2 &:\quad x^{-2} + x^{-1} + 1 + x + \cdots = \frac{x^{-2}}{1 - x}, \\
c_2 \ge -4 &:\quad x^{-4} + x^{-3} + \cdots = \frac{x^{-4}}{1 - x}, \\
-3 \le c_3 \le 6 &:\quad x^{-3} + x^{-2} + \cdots + x^{6} = \frac{x^{-3}(1 - x^{10})}{1 - x}, \\
c_4 \ge 0 &:\quad 1 + x + x^2 + \cdots = \frac{1}{1 - x}.
\end{align*}
相乘得 $f(x) = \dfrac{x^{-2} \cdot x^{-4} \cdot x^{-3}(1 - x^{10})}{(1-x)^4}
= \dfrac{x^{-9}(1 - x^{10})}{(1 - x)^4}$，答案是 $f(x)$ 裡 $x^{23}$ 的係數。 $c_3$ 有 $6 - (-3) + 1 = 10$
個值，所以是 $1 - x^{10}$。

\textbf{換成非負變數}。 令 $y_1 = c_1 + 2$、$y_2 = c_2 + 4$、$y_3 = c_3 + 3$、$y_4 = c_4$，
則 $y_1, y_2, y_4 \ge 0$、$0 \le y_3 \le 9$，方程式變成 $y_1 + y_2 + y_3 + y_4 = 23 + 2 + 4 + 3 = 32$。
對應的生成函數是 $g(x) = \dfrac{1 - x^{10}}{(1 - x)^4} = x^{9} f(x)$，要 $x^{32}$ 的係數，和上面
一樣。 很多教科書要的「生成函數」是這一個形式，兩種寫法都要說清楚要哪一個係數。

\textbf{算出係數}（題目只要生成函數，這段是附帶的）。 $\dfrac{1}{(1-x)^4} = \sum_{k \ge 0}
\binom{k + 3}{3} x^k$，所以
\[
[x^{32}]\, g(x) = \binom{32 + 3}{3} - \binom{22 + 3}{3} = \binom{35}{3} - \binom{25}{3}
= 6545 - 2300 = 4245 .
\]
第二項是扣掉 $y_3 \ge 10$ 的解（$x^{10}$ 那一項把指數往上推了 $10$）。
\end{pastsoln}
